%%%%%%%%%%%%%%%%%%%%%%%%%%%%%%%%%%%%%%%%%%%%%%%%%%%%%%%%%%%%%%%%%
\chapter{Unit 4: Designing Interpretable Predictive Models}
\label{chap:unit4}
%%%%%%%%%%%%%%%%%%%%%%%%%%%%%%%%%%%%%%%%%%%%%%%%%%%%%%%%%%%%%%%%%

\textbf{Total Time: 5 hours} \\
\textbf{Instructional Time: 3.5 hours} \\
\textbf{Assessment Time: 1.5 hours}

This unit introduces the foundations of supervised machine learning
models, with a focus on interpretability and communicating modeling
decisions. Throughout the unit, the \textbf{iBudget Florida Medicaid
Waiver Study} (Scientific Report, Volumes~I and~II, available as
reference documents) serves as the primary worked example. The
iBudget study evaluated the algorithm that Florida uses to allocate
support budgets to over 35,000 individuals with developmental
disabilities, building and comparing eight predictive models using
real administrative data. It illustrates in a high-stakes, real-world
setting every concept this unit covers: baseline model construction,
interpretable feature engineering, multicollinearity diagnostics,
model comparison, and stakeholder reporting. Volume~II, which extends
the predictive results into a resource allocation framework, is offered
as supplemental reading for awareness; students are not expected to
implement its methods. The parallel task throughout this unit is the
2026 Data Challenge, in which students apply the same workflow to NCHS
vital statistics data to project underweight births and infant
mortality by Texas county.

%================================================================
\section{Foundations of Predictive Modeling}
\label{sec:4.1}
%================================================================

\textbf{Time: 1.5 hours (Instructional: 70 minutes, Assessment:
20 minutes)}

\noindent This lesson introduces supervised learning through the
workflow that students will use throughout the rest of the unit:
selecting a prediction target, choosing features, fitting an
interpretable baseline model, and evaluating it on held-out data.
The emphasis is on understanding what the model is doing, not only
on obtaining a good score.

\subsection{Learning Objectives}

\begin{enumerate}
    \item Distinguish supervised learning from other common data
    analysis tasks, and distinguish regression from classification.
    \item Before fitting any model, articulate the prediction problem
    in research terms: name the outcome variable and how it is
    operationalized in the dataset, identify the unit of analysis,
    describe the population represented, and state the intended use
    of the forecast.
    \item Build a linear regression model as an interpretable baseline
    and explain what its coefficients represent.
    \item Practice the workflow for performing a train-test split,
    training a model on training data, evaluating it on held-out test
    data, and understanding how cross-validation works in practice.
    \item Identify common risks in predictive modeling, including
    overfitting and data leakage.
\end{enumerate}

\subsection{Instructional Notes}

The iBudget study opens with exactly the framing exercise described
in Learning Objective~2. Before any model is built, the study
defines its target variable (annualized support cost per enrolled
client), its unit of analysis (individual clients in the Florida
Home and Community-Based Services waiver), and the operational use
of the prediction: the model output becomes the need estimate that
informs each client's budget allocation. This framing step is not
bookkeeping; it determines what a good model means. A model that
predicts well on average but systematically underpredicts costs for
the highest-need clients may be statistically respectable and
operationally harmful at the same time.

For the data challenge, students face an analogous framing decision.
The outcome could be operationalized as the county-level rate of
underweight births (a proportion), as the raw count of underweight
births, or as a binary indicator of whether a county exceeds a
threshold rate. Each operationalization implies a different model
family and a different error cost. Students should make this choice
explicitly and record it before building their baseline.

The iBudget baseline is the frozen 2015 model: a square-root
transformed linear regression with coefficients fixed at their 2015
values, applied to FY2024--25 data without recalibration. This is
a meaningful baseline, not a straw man. It answers the question of
whether the operational model has drifted from the conditions under
which it was calibrated. Students should draw the analogy to their
own baseline choice: what is the simplest defensible model, and
what question does running it answer?

The 2015 iBudget study reported only in-sample metrics, which the
2025 re-evaluation identifies as a critical limitation. This is a
concrete illustration of why the train-test split discipline in
Learning Objective~4 matters: in-sample fit is not evidence of
predictive accuracy, and a model evaluated only on its training data
cannot be trusted to generalize.

\subsection{Assessment Instrument (20 minutes)}

Students will build a baseline linear regression model using the
Vital Statistics dataset. They will document the outcome variable
and explain how it is operationalized in the NCHS records, list the
features used, report one test-set performance metric such as RMSE
or MAE, and write a brief note identifying two ways data leakage or
overfitting could occur in this workflow.

%================================================================
\section{Interpretable Feature Engineering and Pipelines}
\label{sec:4.2}
%================================================================

\textbf{Time: 1 hour (Instructional: 40 minutes, Assessment:
20 minutes)}

\noindent This lesson focuses on interpretable feature engineering:
creating features that better reflect patterns in the data while
preserving interpretability. Students will use plots and domain
knowledge to motivate simple transformations, encodings, and
interactions, and will organize these steps within a small
\verb|sklearn| Pipeline.

\subsection{Learning Objectives}

\begin{enumerate}
    \item Use visualizations to spot patterns in the data that may
    motivate feature engineering, such as nonlinearity, thresholds,
    and group differences.
    \item Understand when to use simple transformations, interaction
    terms, and common encoding strategies such as one-hot encoding.
    \item Build a small preprocessing and modeling Pipeline in
    \verb|sklearn| and explain how the resulting features affect
    model interpretability.
\end{enumerate}

\subsection{Instructional Notes}

The iBudget study's feature engineering process illustrates how
visualizations motivate transformations. Before any model was fit,
the study generated box plots of budget utilization by living
setting, by age group, by race and ethnicity, and by county, and a
scatter plot of budgeted amount against actual spending. Each of
these plots revealed a pattern that motivated a feature decision.
Utilization varied substantially by living setting, which motivated
including it as a predictor. The budget-versus-spending scatter
plot revealed a right-skewed outcome distribution, which motivated
considering log and square-root transformations of the response.
Students can apply the same discipline: plot the outcome against
each candidate feature and let the pattern motivate the engineering
decision.

The iBudget study also illustrates a subtler challenge that
one-hot encoding alone does not resolve: categorical variables whose
levels carry different meanings than their labels suggest. The living
setting variable originally had 22 levels, many of which were
sparsely populated and clinically overlapping. The study consolidated
these to 6 levels through a combination of statistical analysis and
domain judgment. One-hot encoding 22 sparse levels would have
produced an unstable and uninterpretable set of dummy variables;
consolidating first, then encoding, produced a feature that was
both informative and legible. Students should look for analogous
situations in the NCHS data, particularly in variables with many
low-frequency categories.

The \verb|sklearn| Pipeline is the right tool for keeping these
transformations organized and reproducible. A Pipeline ensures that
preprocessing steps applied during training are applied identically
to test data, which prevents a common form of data leakage.

\subsection{Assessment Instrument (20 minutes)}

Students will continue working with the Vital Statistics dataset
and will add two or three engineered features in a small
\verb|sklearn| Pipeline. They will submit the pipeline structure
and a short written justification for each added feature, including
why it may improve model fit or interpretability.

%================================================================
\section{Feature Justification, Multicollinearity, and Model
Diagnostics}
\label{sec:4.3}
%================================================================

\textbf{Time: 1.5 hours (Instructional: 60 minutes, Assessment:
30 minutes)}

\noindent Building upon Section~\ref{sec:4.2}, students will learn
how to decide whether a feature should remain in a model. The
emphasis is on practical diagnostics and justification using
correlations, multicollinearity checks, variance inflation factors,
coefficient stability, and residual plots.

\subsection{Learning Objectives}

\begin{enumerate}
    \item Use simple quantitative tools such as Pearson and Spearman
    correlation, multiple $R^2$, and variance inflation factor to
    identify useful, redundant, or unstable features and to recognize
    multicollinearity.
    \item Use residual plots and side-by-side model comparison to
    assess whether added complexity is justified.
    \item Explain why a feature was kept, transformed, or removed
    based on predictive performance, interpretability, and stability.
    \item Examine model performance separately for meaningful
    subgroups, such as geographic region or demographic category, and
    recognize when a model that performs well overall may perform
    unevenly across the population it serves.
\end{enumerate}

\subsection{Instructional Notes}

Multicollinearity is not a pathology to be feared; it is information
about the feature space that must be understood and managed. The
iBudget study screened candidate features using mutual information
with permutation baseline correction across five consecutive fiscal
years, then checked pairwise redundancy using Spearman rank
correlation for continuous-continuous pairs and bias-corrected
Cram\'er's~$V$ for categorical pairs. VIF screening followed for
the final candidate set. This sequence is exactly the workflow
Learning Objective~1 describes; students can treat Chapter~5 of
the iBudget Volume~I as a methodological reference.

The iBudget study also illustrates why comparing more than two
models is scientifically honest. The study evaluated eight models
under a common protocol: an 80/20 train-test split with 10-fold
cross-validation, all metrics computed on the held-out test set.
The comparison included OLS re-estimation, a Generalized Linear
Model with a Gamma distribution (which handles the positive,
right-skewed distribution of support costs more appropriately than
OLS), robust regression using Huber M-estimation, ridge regression
to manage multicollinearity among correlated assessment items, and
a random forest as the accuracy ceiling. Students do not need to
implement all of these, but they should understand why a family
of models exists: each alternative addresses a specific limitation
of the OLS baseline that can be diagnosed from residual plots and
distributional checks. A fitted-vs-residual plot that shows a
funnel pattern suggests heteroscedasticity; a GLM or weighted least
squares may be warranted. A pattern of large positive residuals for
the highest-cost observations suggests the outcome distribution
needs a transformation or a different family entirely.

Learning Objective~4 reflects a lesson that emerges clearly from
the iBudget subgroup analysis. A model that predicts overall support
costs accurately can still systematically underpredict costs for the
highest-need clients, whose budgets have the largest consequences
for service access. In the data challenge context, a county-level
model that performs well on average across Texas may perform poorly
on rural or border counties, which are precisely the populations
where targeted public health intervention is most important. Subgroup
performance is not optional reporting; it is part of what the model
works means.

\subsection{Assessment Instrument (30 minutes)}

Students will compare two candidate models, such as a baseline
model and an expanded feature-engineered model. They will decide
which model they would keep and justify the decision using at least
three pieces of evidence drawn from performance metrics,
multicollinearity checks, coefficient stability, or residual
diagnostics. The justification should also include one observation
about whether the preferred model's performance is consistent across
at least one meaningful subgroup in the data.

%================================================================
\section{Model Reporting and Stakeholder Communication}
\label{sec:4.4}
%================================================================

\textbf{Time: 1 hour (Instructional: 30 minutes, Assessment:
30 minutes)}

\noindent This lesson focuses on communicating model performance
and design decisions to external stakeholders. The goal is to turn
the analytic work from Section~\ref{sec:4.3} into a clear written
summary that reports model results responsibly and in a form that
aligns with the broader reporting expectations introduced earlier
in the curriculum.

\subsection{Learning Objectives}

\begin{enumerate}
    \item Understand how to report and compare interpretable
    predictive models using appropriate regression metrics such as
    MSE, RMSE, and MAE.
    \item Summarize model design choices, strengths, and limitations
    in language appropriate for a non-technical audience.
    \item Practice connecting model reporting to broader ideas of
    transparency and reproducibility, including the spirit of
    TRIPOD-style reporting.
\end{enumerate}

\subsection{Instructional Notes}

The iBudget study's reporting structure is a professional example of
the TRIPOD spirit applied to an administrative prediction problem.
Each model in the comparison follows an identical chapter template:
specification, overall performance metrics, accuracy bands,
cross-validation results, subgroup analysis, variance and stability
metrics, and implementation considerations. This structure means
that any reader can navigate to any modeling decision and find it
documented. Nothing is left to be inferred from code. Students do
not need to produce a report of that scope, but they can use it as
a reference for what complete documentation looks like and what
questions a careful reader will ask.

The report produced in this section will be carried forward. Unit~6
asks whether estimates are stable across data sources and time
periods, and a model report that does not specify its outcome
operationalization or validation strategy cannot be meaningfully
compared to estimates from other sources. Unit~7 subjects the report
to LLM ensemble critique, and a report with gaps will produce
ensemble outputs that are difficult to verify. Writing the report
well here has direct consequences for the quality of work in both
subsequent units.

\subsection{Assessment Instrument (30 minutes)}

Students will write a short report for a stakeholder describing
their final model. The report should identify the prediction target,
summarize the features included, report one or two performance
metrics, and explain at least one important limitation or caveat in
plain language. The report should also include one sentence
connecting the model to the data challenge: what specific projection
will this model contribute to the final analysis, and for which
population or geography?

% Required packages (add to your preamble if not already there):
% \usepackage{longtable}
% \usepackage{booktabs}
% \usepackage{ragged2e}
% \usepackage{array}
% \usepackage{pdflscape}

\begin{landscape}

    \section{Evaluation Rubric}

    {\footnotesize

    \begin{longtable}{%
        >{\RaggedRight\arraybackslash}p{2.5cm}   % Component
        >{\RaggedRight\arraybackslash}p{4.5cm}   % Excellent
        >{\RaggedRight\arraybackslash}p{4.2cm}   % Good
        >{\RaggedRight\arraybackslash}p{4.2cm}   % Satisfactory
        >{\RaggedRight\arraybackslash}p{4.5cm}}  % Needs Improvement

    \toprule
    \textbf{Component} &
    \textbf{Excellent (90--100\%)} &
    \textbf{Good (80--89\%)} &
    \textbf{Satisfactory (70--79\%)} &
    \textbf{Needs Improvement ($<$70\%)} \\
    \midrule
    \endfirsthead

    \toprule
    \textbf{Component} &
    \textbf{Excellent (90--100\%)} &
    \textbf{Good (80--89\%)} &
    \textbf{Satisfactory (70--79\%)} &
    \textbf{Needs Improvement ($<$70\%)} \\
    \midrule
    \endhead

    \midrule
    \multicolumn{5}{r}{\textit{Continued on next page}} \\
    \endfoot

    \bottomrule
    \endlastfoot

    % Row 1
    \textbf{1.\ Foundations of Predictive Modeling}
    \newline\textbf{(25\%)}
    &
    Clearly distinguishes supervised vs.\ unsupervised, regression
    vs.\ classification; correctly performs a train-test split, trains an
    interpretable baseline model, evaluates it on held-out data, and explains
    how cross-validation works; articulates the prediction problem in research
    terms before fitting; demonstrates strong understanding of overfitting
    and data leakage.
    &
    Good workflow with minor errors in reasoning or terminology;
    prediction framing present but partially developed;
    overfitting/data leakage described but not deeply.
    &
    Workflow partially correct; inconsistent understanding of splits,
    evaluation metrics, or model fitting; prediction framing absent
    or generic.
    &
    Incorrect or incomplete workflow; poor understanding of model evaluation,
    overfitting, or data leakage; no research framing of the prediction
    problem.
    \\
    \midrule

    % Row 2
    \textbf{2.\ Interpretable Feature Engineering \& Pipelines}
    \newline\textbf{(20\%)}
    &
    Insightfully uses visualizations to identify meaningful patterns;
    appropriately applies simple transformations, encoding strategies, and
    interaction terms; builds a clear \verb|sklearn| Pipeline; thoroughly documents
    the rationale for each feature added.
    &
    Good feature construction with minor justification gaps; Pipeline mostly
    correct; documentation adequate.
    &
    Basic feature creation with limited insight; Pipeline incomplete or
    justification superficial.
    &
    Poor or incorrect feature construction; missing Pipeline; unclear or
    unjustified decisions.
    \\
    \midrule

    % Row 3
    \textbf{3.\ Feature Justification, Multicollinearity \& Diagnostics}
    \newline\textbf{(25\%)}
    &
    Correctly applies feature justification tools such as correlation checks,
    multiple $R^2$, variance inflation factor, and residual diagnostics;
    correctly identifies multicollinearity concerns; compares models
    thoughtfully and gives a clear, evidence-based rationale for keeping the
    preferred model; includes at least one subgroup performance observation
    with genuine analytical content.
    &
    Most diagnostics applied correctly; model comparison reasonable but with
    minor gaps in interpretation or justification; subgroup observation
    present with adequate reasoning.
    &
    Partial or inconsistent use of diagnostics; model choice weakly
    justified; subgroup observation absent or superficial.
    &
    Incorrect or missing diagnostic evidence; model choice unsupported or
    inconsistent with results; no subgroup analysis.
    \\
    \midrule

    % Row 4
    \textbf{4.\ Model Reporting \& Stakeholder Communication}
    \newline\textbf{(20\%)}
    &
    Accurately computes and interprets regression metrics such as MSE, RMSE,
    and MAE; report is clear, well-structured, and communicates model design
    decisions, performance, and limitations appropriately for the intended
    audience; includes a specific, meaningful connection to the data challenge.
    &
    Metrics computed correctly with minor errors; report mostly clear and
    appropriately framed; data challenge connection present but general.
    &
    Metrics included but inconsistently computed or interpreted; explanation
    limited or too technical; data challenge connection absent or formulaic.
    &
    Metrics missing or incorrect; report unclear, incomplete, or not aligned
    with results; no connection to the data challenge.
    \\
    \midrule

    % Row 5
    \textbf{5.\ Overall Professionalism, Clarity, and Documentation}
    \newline\textbf{(10\%)}
    &
    Writing polished and well-organized; notebooks/scripts clean,
    reproducible, and well-commented; outputs easy to interpret.
    &
    Mostly clear with minor structural or formatting issues; code readable.
    &
    Understandable but inconsistent structure or documentation.
    &
    Poorly organized, unclear writing, missing documentation, or sloppy code.
    \\

    \end{longtable}

    } % end \footnotesize

\end{landscape}